\section{Parte 3: Product Backlog priorizado}

Ordenamos las historias por prioridad usando MoSCoW. Las de tipo \textbf{Must} son las que sí o sí necesitamos para que el sistema funcione, las \textbf{Should} son importantes pero no bloquean, y las \textbf{Could} son mejoras que podemos dejar para después.

\begin{table}[htbp]
\centering
\small
\renewcommand{\arraystretch}{1.15}
\begin{tabularx}{\textwidth}{c c c X}
\toprule
\textbf{\#} & \textbf{ID} & \textbf{Prioridad} & \textbf{Historia de usuario} \\
\midrule
1 & HU-10 & Must & Asignar roles a los usuarios para controlar el acceso. \\
\midrule
2 & HU-06 & Must & Crear capacitaciones con fechas, cupos y docente. \\
\midrule
3 & HU-01 & Must & Inscripción de participantes con DNI. \\
\midrule
4 & HU-03 & Must & Pasar asistencia con código QR. \\
\midrule
5 & HU-04 & Must & Registrar notas y calcular aprobación. \\
\midrule
6 & HU-05 & Must & Generar certificado PDF con código QR. \\
\midrule
7 & HU-07 & Should & Gestionar varias instituciones en el sistema. \\
\midrule
8 & HU-08 & Should & Verificación pública de certificados. \\
\midrule
9 & HU-09 & Should & Reportes para auditorías. \\
\midrule
10 & HU-02 & Should & Confirmación de inscripción por correo. \\
\midrule
11 & HU-12 & Could & Catálogo de cursos disponibles. \\
\midrule
12 & HU-11 & Could & Asistencia offline con sincronización. \\
\bottomrule
\end{tabularx}
\caption{Product Backlog priorizado del SIGC-CUSCO.}
\label{tab:backlog}
\end{table}
